% P11 paper module: R12 positive resolution of the logarithmic complement gate.

\subsection{Positive Sobolev bootstrap for the explicit complement}\label{subsec:o3k-log-complement}

The previous subsection deliberately stopped at the operator-domain gain
\eqref{eq:o3j-domain-gain}.  At fixed terminal horizon, however, the source cutoff has
additional structure: after resolving the finite martingale coordinates it is a finite
family of interval cutoffs and translations.  This permits a small positive Sobolev
bootstrap.  Since any positive Sobolev order dominates every finite logarithmic order,
this closes the logarithmic complement gate.

For $s\in\mathbb R$ let
\[
 X_T^s:=\{F\in H^s(\mathbb R):\operatorname{supp}F\subset[-T,T]\}
\]
with the inherited $H^s$ norm.  We use the following standard cutoff fact.

\begin{lemma}[Small Sobolev cutoff range]\label{lem:o3k-cutoff}
For every bounded interval $I\subset\mathbb R$, multiplication by $1_I$ is bounded on
$H^s(\mathbb R)$ whenever $|s|<1/2$.  More generally, for fixed finite $a\ge0$ it is
bounded on the weighted Fourier space
\[
 \mathcal H^{s,a}
 :=\left\{F:\int_{\mathbb R}
 \langle\xi\rangle^{2s}[\log(2+|\xi|)]^{2a}
 |\widehat F(\xi)|^2\,d\xi<\infty\right\}
\]
whenever $|s|<1/2$.
\end{lemma}

\begin{proof}
The first statement is the classical fractional-Sobolev interval-multiplier theorem.
For the logarithmically refined statement put
\[
 w_{s,a}(\xi):=\langle\xi\rangle^{2s}[\log(2+|\xi|)]^{2a}.
\]
For $|s|<1/2$, $w_{s,a}$ is an $A_2$ weight: a direct interval-average estimate reduces
at large scale to the power range $-1<2s<1$, while the logarithmic factor is slowly
varying and does not alter the open $A_2$ range.  Multiplication by a half-line in
physical space becomes, after Fourier transform, a fixed linear combination of the
identity and the Hilbert transform.  The weighted Hilbert-transform theorem on
$L^2(w_{s,a})$ therefore gives boundedness for half-line cutoffs; differences of two
translated half-lines give bounded intervals.
\end{proof}

\begin{lemma}[Fixed-window Schur Sobolev window]\label{lem:o3k-schur-window}
For every fixed $T>0$ there exists $\varepsilon_T>0$ such that
\[
 B_T=(I+R_T^*R_T)^{-1}:X_T^s\longrightarrow X_T^s
\]
is bounded for all $|s|<\varepsilon_T$.  Consequently $H_T$, $H_T^*$ and
$\Sigma_T=H_TB_TH_T^*$ are bounded on $X_T^s$ after possibly shrinking
$\varepsilon_T$.
\end{lemma}

\begin{proof}
At fixed $T$ the sum defining $R_T$ is effectively finite.  By
\eqref{eq:projected-mark}, each martingale coordinate contains a scalar factor
\[
 1_{\{j<J_{p,T}(u)\}},
\]
which, up to endpoints, is the indicator of
\[
 \left\{|u|<T-\frac{j+1}{2}\log p\right\}.
\]
Thus, after resolving the finite martingale coordinates, $R_T$ is a finite sum of
compositions of fixed translations, interval cutoffs, zero extension/restriction and
fixed finite-dimensional coefficient maps.  Lemma~\ref{lem:o3k-cutoff} therefore gives
bounded maps
\[
 R_T:X_T^s\longrightarrow Z_T^s
 \qquad(|s|<1/2)
\]
on the corresponding finite vector-valued Sobolev target.  Applying the same statement
at exponent $-s$ and dualizing gives
\[
 R_T^*:Z_T^s\longrightarrow X_T^s.
\]
Hence
\[
 A_T^{\rm rest}:=I+R_T^*R_T
\]
is bounded on $X_T^s$ for $|s|<1/2$.

At $s=0$, $A_T^{\rm rest}\ge I$, so it is an isomorphism on $X_T^0=L^2(-T,T)$.
The spaces $X_T^s$ form a complex interpolation scale for $|s|<1/2$: they are the
ranges of the common interval-cutoff projection, which is bounded at the endpoint
exponents by Lemma~\ref{lem:o3k-cutoff}.  The local stability theorem for invertible
operators on complex interpolation scales (the \v Sne\u\i berg theorem) now yields an
$\varepsilon_T>0$ such that $A_T^{\rm rest}$ remains invertible on $X_T^s$ whenever
$|s|<\varepsilon_T$.  This is the asserted boundedness of $B_T$.

The hub $H_T$ is itself a finite sum of fixed translation differences and interval
extension/restriction, hence is bounded on the same small Sobolev scale.  Boundedness of
$H_T^*$ follows by using the negative exponents and duality, and the assertion for
$\Sigma_T$ follows by composition.
\end{proof}

\begin{proposition}[Positive Sobolev regularity of the O3j forcing]\label{prop:o3k-forcing}
For the fixed triple $0<R<S<T_0$ and smooth annulus vector $h$ of
Subsection~\ref{subsec:o3j-reconciled}, there exists $s_1>0$ such that
\[
 \boxed{E_Rr_h\in H^{s_1}(\mathbb R).}
\tag{OJ.12}\label{eq:o3k-forcing-sobolev}
\]
\end{proposition}

\begin{proof}
Put $\phi=E_{S,T_0}h\in C_c^\infty((-T_0,T_0))$.  Choose
$0<s_1<\min\{\varepsilon_{T_0},1/2\}$ in
Lemma~\ref{lem:o3k-schur-window}.  Since $\phi$ is smooth and supported strictly inside
$(-T_0,T_0)$,
\[
 H_{T_0}^*\phi\in X_{T_0}^{s_1}.
\]
Therefore
\[
 \Sigma_{T_0}\phi
 =H_{T_0}B_{T_0}H_{T_0}^*\phi
 \in X_{T_0}^{s_1}.
\]
Restriction to $(-R,R)$ and zero extension back to the line preserve $H^{s_1}$, after
shrinking $s_1$ if necessary, so
\[
 E_Rr_{\sigma,h}\in H^{s_1}(\mathbb R).
\]
On the Gamma side, Lemma~\ref{lem:gamma-smooth-action} gives
$r_{\Gamma,h}=P_R\mathcal G_\phi$ with
$\mathcal G_\phi\in H^N(\mathbb R)$ for every finite $N$.  Lemma~\ref{lem:o3k-cutoff}
therefore gives $E_Rr_{\Gamma,h}\in H^{s_1}$ for every sufficiently small
$s_1<1/2$.  Adding the two terms proves~\eqref{eq:o3k-forcing-sobolev}.
\end{proof}

We next bootstrap through the genuine finite-window form operator, rather than replacing
it by a full-line Fourier multiplier.  Put $m=m_\Gamma$ and for $|s|<1/2$ define
\[
 \mathcal V^s
 :=\left\{F:\int_{\mathbb R}
 \langle\xi\rangle^{2s}m(\xi)|\widehat F(\xi)|^2\,d\xi<\infty\right\},
\qquad
 V_R^s:=\{F\in\mathcal V^s:\operatorname{supp}F\subset[-R,R]\}.
\tag{OJ.13}\label{eq:o3k-Vs}
\]
Since $m(\xi)\asymp\log(2+|\xi|)$, Lemma~\ref{lem:o3k-cutoff} applies to these weights.
The $V_R^s$ are therefore complemented support subspaces of the full weighted Fourier
scale and satisfy the complex interpolation law
\[
 [V_R^{-\delta},V_R^{\delta}]_\theta
 =V_R^{(2\theta-1)\delta}
\tag{OJ.14}\label{eq:o3k-interpolation}
\]
for every sufficiently small fixed $\delta>0$.  Let
\[
 Y_R^s:=(V_R^{-s})^*.
\]

\begin{theorem}[Positive Sobolev regularity of the explicit complement]\label{thm:o3k-complement}
For the explicit vector
\[
 g_h=h-J_{R,S}u_h
\]
there exists a number
\[
 0<s_*=s_*(R,S,T_0,h)<1/2
\]
such that
\[
 \boxed{E_Ru_h\in H^{s_*}(\mathbb R),\qquad
 E_Sg_h\in H^{s_*}(\mathbb R).}
\tag{OJ.15}\label{eq:o3k-positive-sobolev}
\]
Consequently
\[
 \boxed{
 E_Sg_h\in\bigcap_{\alpha<\infty}\HG^\alpha(\mathbb R).
 }
\tag{OJ.16}\label{eq:o3k-all-log}
\]
In particular, if $m_h$ denotes the first nonzero integral-jet order of $g_h$, then
\[
 \boxed{E_Sg_h\in\HG^{m_h+3/2}(\mathbb R).}
\tag{OJ.17}\label{eq:o3k-log-gate}
\]
\end{theorem}

\begin{proof}
For $u\in V_R^s$ and $v\in V_R^{-s}$ define
\[
 \langle\mathcal A_Ru,v\rangle
 :=\mathfrak c_{\Gamma,R}[u,v]
 +\langle\Sigma_R^{[T_0]}u,v\rangle.
\tag{OJ.18}\label{eq:o3k-form-operator}
\]
The Gamma part is bounded by Fourier-space Cauchy--Schwarz:
\[
 |\mathfrak c_{\Gamma,R}[u,v]|
 \le \|u\|_{V_R^s}\|v\|_{V_R^{-s}}.
\]
By Lemma~\ref{lem:o3k-schur-window}, zero extension/restriction between the fixed nested
intervals and the factorization
\[
 \Sigma_{T_0}=H_{T_0}B_{T_0}H_{T_0}^*
\]
show, after shrinking $\delta>0$, that
$\Sigma_R^{[T_0]}$ is bounded on $H^s(-R,R)$ for $|s|\le\delta$.  Since $m\ge1$,
\[
 V_R^s\hookrightarrow H^s,
 \qquad V_R^{-s}\hookrightarrow H^{-s},
\]
so the $H^s$--$H^{-s}$ dual pairing gives
\[
 |\langle\Sigma_R^{[T_0]}u,v\rangle|
 \le C_s\|u\|_{V_R^s}\|v\|_{V_R^{-s}}.
\]
Thus
\[
 \mathcal A_R:V_R^s\longrightarrow Y_R^s
\]
is bounded throughout a symmetric interpolation interval around zero.

At $s=0$, $V_R^0$ is precisely the Gamma form space under zero extension, and
\[
 \langle\mathcal A_Ru,u\rangle
 =\mathfrak c_{\Gamma,R}[u]+\langle\Sigma_R^{[T_0]}u,u\rangle
 \ge \|u\|_{V_R^0}^2.
\]
Lax--Milgram therefore gives an isomorphism
\[
 \mathcal A_R:V_R^0\xrightarrow{\sim}Y_R^0.
\]
Apply the \v Sne\u\i berg local-stability theorem to the interpolation scales
\eqref{eq:o3k-interpolation} and their dual scales.  There exists $\varepsilon_1>0$ such
that
\[
 \mathcal A_R:V_R^s\xrightarrow{\sim}Y_R^s
 \qquad(|s|<\varepsilon_1).
\tag{OJ.19}\label{eq:o3k-form-invertible}
\]

By Proposition~\ref{prop:o3k-forcing}, $E_Rr_h\in H^{s_1}$ for some $s_1>0$.
Choose
\[
 0<s_*<\min\{s_1,\varepsilon_1\}.
\]
Then $r_h$ defines an element of $Y_R^{s_*}$ by the usual
$H^{s_*}$--$H^{-s_*}$ pairing.  Equation~\eqref{eq:o3k-form-invertible} therefore has a
unique solution $u_*\in V_R^{s_*}$.  Since $V_R^{s_*}\subset V_R^0$ and
$V_R^0\subset V_R^{-s_*}$, testing the $s_*$ equation against every vector in $V_R^0$
shows that $u_*$ solves the original weak equation~\eqref{eq:o3j-forcing-equation}.
Uniqueness at $s=0$ gives $u_*=u_h$.  Hence
\[
 E_Ru_h\in V_R^{s_*}\subset H^{s_*}(\mathbb R).
\]
Because
\[
 E_Sg_h=E_Sh-E_Ru_h
\]
and $E_Sh$ is Schwartz, the first assertion of~\eqref{eq:o3k-positive-sobolev} follows.
Finally, for every finite $\alpha$,
\[
 [\log(2+|\xi|)]^{2\alpha}
 \le C_{\alpha,s_*}\langle\xi\rangle^{2s_*},
\]
which proves~\eqref{eq:o3k-all-log} and hence~\eqref{eq:o3k-log-gate}.
\end{proof}

\begin{remark}[R12 firewall]\label{rem:o3k-firewall}
The exponent $s_*>0$ is a fixed-window quantity and may depend on $R,S,T_0$ and on the
chosen complement vector; no lower bound uniform in terminal horizons is asserted.
Theorem~\ref{thm:o3k-complement} closes the logarithmic regularity gate used by the O3
second-moment diagnostic.  It does not control the polar unitary gauges of
Remark~\ref{rem:polar-gauge}, does not prove $K_{R,S}^{T,U}\to I$, and does not prove or
disprove strong terminal transport.
\end{remark}
